%% kapa: exam=1 round=37 year=2026 period=2 form=A subject=appraisal_law subject_ko=감정평가관계법규 seq=13 number=13
%% source: Q-Net 감정평가사 기출문제 (cst00309) artlSeq=5258924 file=raw/1008/5258924/2026년 제37회 감정평가사 1차 2교시.pdf
%% conversion: kordoc 4.24.51 -> tools/gen1.py (auto); review=proofread-sonnet; flags=none
\begin{question}{13}
국토의 계획 및 이용에 관한 법령상 중앙도시계획위원회에 관한 설명으로 옳은 것은?
\begin{choices}
  \item 위원장과 부위원장이 모두 부득이한 사유로 그 직무를 수행하지 못할 때에는 위원장이 미리 지명한 위원이 그 직무를 대행한다.
  \item 위원회는 위원장ㆍ부위원장 각 1명을 포함한 15명 이상 20명 이하의 위원으로 구성한다.
  \item 위원장이 각 분과위원회의 위원장을 지명한다.
  \item 공무원이 아닌 위원의 수는 5명 이상으로 하고, 그 임기는 1년으로 한다.
  \item 보궐위원의 임기는 1년으로 한다.
\end{choices}
\end{question}
\answer{1}
